\section{The metric can remove irrational-time weak modes}
\label{sec:full-metric-holonomy}

The irrational-time obstruction in
Proposition~\ref{prop:both-ray-holonomy}
concerns observations of the orbit metric alone.
The remaining spatial metric coefficient gives
a scalar reading of the otherwise hidden
transported velocity. This reading has a
different principal symbol from the two orbit
coordinates.

Use the notation of Section~\ref{sec:both-velocities}.
In addition to both orbit fields at $c$,
measure the metric gradient $\lambda_\theta(c)$
in $H^s$, equivalently the mean-zero metric
potential in $H^{s+1}$. Use the initial orbit
fields and a metric anchor at $a$.
After the full orbit observation is used,
the principal metric reading is proportional to
\[
 B(\theta)\cdot H(\theta),\qquad
 B(\theta)=A_-^{-1}(c,\theta)U_\theta(c,\theta).
\]
In fact the derivative of the momentum constraint,
after taking a spatial primitive, gives to principal order
\[
 -\Big(\langle A_-H,U_t+U_\theta\rangle
       -\langle A_+\mathcal T H,U_t-U_\theta\rangle\Big).
\]
The contribution involving $U_t-U_\theta$ paired
with $A_-H-A_+\mathcal T H$ is already observed.
The remaining contribution is $-2 B\cdot H$.
The common scale of this expression is $\sqrt{ac}$.
Differentiating the transported primitive gives
this expression by the momentum equation, whose
linearization is
$-2t(\langle D_t V,U_\theta\rangle+
\langle U_t,D_\theta V\rangle)$.
The differentiated transport coefficients and
the forced error in \eqref{eq:transport-leading}
give terms bounded from $H^{s-1}$ to $H^{s+1}$
after the primitive. The mean is a finite-rank term.

\begin{theorem}[The metric projection along successive rays]
\label{th:full-metric-gramian}
Let $d/\pi$ be irrational and let
$B_j(\theta)=(\mathcal U_d^j B)(\theta)$.
The linearized observation of the terminal
orbit fields and metric gradient is bounded below
in the $H^s$ velocity norm up to an
$H^{s-1}$ remainder if and only if
for some finite $J$ and $b>0$,
\begin{equation}
 \sum_{j=0}^J B_j(\theta)B_j(\theta)^\mathsf{T}\ \ge bI
          \quad\text{for every }\theta .
 \label{eq:metric-gramian}
\end{equation}
The sufficient direction holds at rational delays
as well. Under \eqref{eq:metric-gramian},
finitely many early orbit coefficients remove
any finite-dimensional kernel, giving a local
nonlinear Lipschitz inverse for both velocities.
For sufficiently short irrational delays, the
criterion holds if $U_\theta(a)$ and
$D_\theta U_\theta(a)$ are independent at
some angle of the initial orbit curve. In particular
this applies at the common initial node of the
Ernst pair, where the orbit curve is a nonconstant
horocycle arc in the upper half-plane.
\end{theorem}

\begin{proof}
Write $\mathcal U=\mathcal U_d$.
Orthogonality of the transport gives the identity
\[
 B_j\cdot H
    =\mathcal T^j(B\cdot H)
         +B_j\cdot(H-\mathcal U^j H).
\]
The second term is controlled by telescoping
$H-\mathcal U H$ and boundedness of $B_j$.
Condition \eqref{eq:metric-gramian} then controls
$\|H\|_{L^2}$ from the two readings.
Applying the same argument to the top spatial
derivative gives the Sobolev estimate up to
one lower order. The smoothing terms in the
linearized wave and metric maps give the
claimed $H^{s-1}$ remainder.
Reading the early low orbit modes and taking
the nonlinear neighborhood
uses the same proof as
Theorem~\ref{th:fixed-rank-two}, with two
components of initial velocity in the energy.

Conversely, if \eqref{eq:metric-gramian}
fails for every $J$, for arbitrarily long finite
orbit segments there is a unit vector at one end
orthogonal to all transported $B$ readings on
that segment. Transport this vector back to the
first angle of the segment. Place an $L^2$-unit
smooth envelope $h_0$ around that angle, pointing
in the transported vector direction. The supports
of $\mathcal U^j h_0$ for $0\le j\le J$ can be
made disjoint by irrationality.
Their $B$ readings are uniformly as small as
desired by continuity and the orthogonality at
the centers, after choosing the envelope narrow
enough for this finite segment.
The normalized sum
$(J+1)^{-1/2}\sum_{j=0}^J\mathcal U^j h_0$
has unit norm and its difference under $I-\mathcal U$
has norm at most $2/\sqrt{J+1}$.
Multiplication by a high Fourier carrier with
$e^{-2imd}$ arbitrarily near $1$ preserves these
bounds up to an arbitrarily small error.
Taking the carriers sufficiently high yields a
weakly null sequence in $L^2$.
Conjugating to $H^{s+1}$ and imposing the
finite mean and momentum restrictions introduce
only compact or finite-rank corrections, which
vanish in this sequence. The claimed lower
estimate must therefore fail.

To see the short-delay criterion, compare
$B_0$ and $B_1$ in the terminal frame at
$(c,\theta)$ using $A_-(c,\theta)$.
They become $U_\theta(c,\theta)$ and
$U_\theta(c,\theta-2d)$ transported to $(c,\theta)$
across the two rays with common initial point
$(a,\theta-d)$. The first-order difference between
these two vectors is $-2dD_\theta U_\theta(a,\theta)$.
Indeed the spatial displacement is $2d$ and the
first-order temporal contributions of the reversed
and forward transports cancel. Thus the two
vectors are independent on an open arc for
sufficiently small $d$. By irrationality every
rotation orbit enters that arc and finitely many
translates cover the circle. Orthogonal transport
preserves independence, giving
\eqref{eq:metric-gramian} with finitely many
readings in the principal estimate.
\end{proof}

\paragraph{The complete high-frequency test for two full metrics.}
If the initial metric profile is supplied, its gradient
also reads the initial velocity through
$\delta\lambda_\theta(a)=-2a\langle f,U_\theta(a)\rangle$.
After taking the spatial primitive and shifting,
this adds the principal row
$B^{(a)}(\theta)=U_\theta(a,\theta+d)$ in the initial
orthonormal frame of $H$.
The extra differentiated coefficient term is
one order smoother. Denote the available rows by
$B^{(\ell)}$, with $\ell$ running over the terminal
metric row and the initial row if supplied.
For irrational $d/\pi$ the same proof gives the
necessary and sufficient test up to compact
remainders with the sum in
\eqref{eq:metric-gramian} taken over these rows.
Their joint mean readings are finite rank.

For rational $d/\pi=p/q$ the corresponding necessary
and sufficient condition up to compact remainders is
uniform positive definiteness of
\begin{equation}
 (I-\mathcal H_d)^\mathsf{T}(I-\mathcal H_d)
 +\sum_{\ell}\sum_{j=0}^{q-1}
   (\mathcal U_d^jB^{(\ell)})
   (\mathcal U_d^jB^{(\ell)})^\mathsf{T}.
 \label{eq:rational-full-metric}
\end{equation}
Indeed $(I-\mathcal H_d)H=H-\mathcal U_d^qH$
is bounded by $q\|(I-\mathcal U_d)H\|$.
Each projected row is controlled by the telescoping
identity in the proof above. Conversely, if the
displayed matrix is singular at an angle, choose
a unit vector fixed by its closed transport and
orthogonal to every projected row.
Use normalized narrow envelopes about its finite
rotation orbit with the transported vector
directions. Choose their relative signs by the
transport and sum them. The residual in
$I-\mathcal U_d$ is small because the last
transport returns the initial direction at the
center of the envelope. All metric readings are
small by orthogonality at the corresponding centers.
Multiplying by arbitrarily high carriers at
multiples of $q$ commutes with the shift and
gives a weakly null sequence. The compact
lower-order corrections then vanish, proving
necessity. This test includes pure closed-ray
holonomy when no metric row is supplied, and
can still be positive at identity holonomy when
the projected metric directions span the plane.

\begin{proposition}[Two slices at irrational separation]
\label{prop:irrational-metric}
For every positive initial radius in the circular-loop
family with zero initial radial velocity, all sufficiently
short positive delays give a local energy-norm
Lipschitz inverse for both velocities from the two
induced metrics. The delays may be rational or irrational.
One may replace the initial metric by the orbit
part and use a single metric anchor. This stability
persists in a neighborhood of the Cauchy data
and on an open interval of nearby known observation
delays.
\end{proposition}

\begin{proof}
In the rotating orthonormal frame of the circular
loop, $B=A_-^{-1}U_\theta$ is constant in $\theta$
and is nonzero when the final radius is positive.
The matrix $W_d$ is also constant in this frame.
Choose the delay short enough that its rotation
angle $-2k_0\int_a^c\cosh R(t)\,dt$ is strictly
between $-\pi$ and zero.
Then $B$ and $W_d B$ are independent and
\eqref{eq:metric-gramian} holds with $J=1$.
It remains to rule out an exact kernel. The linearized
equations decompose in Fourier modes of the
rotating frame. If the terminal orbit and
metric readings vanish, the terminal orbit variation
vanishes and the gradient constraint gives
$\langle D_t V(c),U_\theta(c)\rangle=0$.
Thus for each Fourier mode the terminal derivative
is radial. Evolve the mode backward from zero
terminal position and radial derivative $g e_R$.
In the reversed time variable, after multiplying by
$\sqrt{t/c}$ the exact linearized operator differs
from the frozen-connection operator by a potential
with norm at most $D_0+D_1|k|$, where
\[
 D_0=\sup_{[a,c]}\left(
 |\Omega^2-\Omega_0^2|+\|K\|+\frac{1}{4t^2}\right),
 \qquad
 D_1=2\sup_{[a,c]}|\Omega-\Omega_0|.
\]
Here $\Omega(t)=k_0\cosh R(t)$, $\Omega_0=\Omega(c)$,
and the curvature potential $K$ in the polar
frame has diagonal entries $(k_0\sinh R)^2,-R_t^2$.
Indeed the principal spatial operator on mode $k$
is the square of $ikI+\Omega J$, with $J$ the
rotation generator. The time rescaling removes
$t^{-1}\partial_t$ and adds $1/(4t^2)$.

For the frozen operator the two circular
polarizations have frequencies $k\pm\Omega_0$.
The radial vector has equal projection magnitude
onto them. Put $q_k=\max(d^{-1},|k|)$
and assume $\Omega_0d<\pi/4$ and $\Omega_0d<1/2$.
The free reverse transfer from zero position
is diagonal with entries $\sin((k\pm\Omega_0)r)/
(k\pm\Omega_0)$, using the continuous value at
zero. Its norm times $q_k$ is at most $2$ for
$0\le r\le d$. At $r=d$ the radial input has
output norm times $q_k$ at least
\[
 \frac{\sin(\Omega_0d)}{1+\Omega_0 d}|g|.
\]
This follows by bounding
$|k\pm\Omega_0|\le(1+\Omega_0d)q_k$ and using the
sum of the two squared sines. The weighted
Duhamel formula bounds the correction by
$2\big(\exp(2d(dD_0+D_1))-1\big)|g|$.
The correction is $O(d^2)$ since $D_0$ is bounded
and $D_1=O(d)$, while the free lower bound is of
order $d$. Thus all sufficiently short delays
force nonzero displacement at $a$ for any nonzero
radial terminal derivative, uniformly over $k$.
This rules out the kernel and also gives a direct
lower estimate. For nonzero terminal readings,
the momentum constraint controls the tangential
terminal derivative by the gradient reading and
$q_k$ times the terminal orbit displacement.
Separating that tangential derivative and the measured
displacement in the reverse problem, the same
Duhamel estimate bounds their contribution at $a$
in the $q_k$-weighted position norm. Since the
position at $a$ is prescribed, the radial lower
bound then controls the remaining terminal
derivative by these readings. The regular-time
energy estimate transfers this bound to the initial
velocity. For example $a=1.6$,
$k_0=4$, $R(a)=0.32$, $R_t(a)=0$ allow every
$0<d\le0.02$. The decreasing radial energy
gives $|R|\le0.32$, $|R_t|<1.31$ and $R>0.29$ on
the interval. In particular $4\le\Omega_0<4.21$,
$D_0<2.11$ and $D_1<3.44d$. These bounds
make the correction less than $0.45d$
and the displayed free gain greater than $3.68d$
over that interval.

The compact-remainder bound for the enlarged
metric observation and injectivity now give the
energy-norm lower bound. The nonlinear estimate
again follows from continuity of the Cauchy derivative.

The positive lower bound in the finite sum
\eqref{eq:metric-gramian} persists continuously
in the delay, along with the bounded lower-order
remainder. Translation is strongly continuous on
Sobolev spaces even though it need not be continuous
in operator norm. If the full lower bound failed
arbitrarily close to an injective delay $d_0$, take
unit $H^s$ velocities $f_n$ at $d_n\to d_0$
with terminal observations tending to zero.
After a subsequence $f_n$ converges in $H^{s-1}$
and weakly in $H^s$ to $f$.
The regular-time energy bound and strong
continuity in that weaker norm imply that $f$
is invisible at $d_0$. Injectivity gives $f=0$,
and the uniform lower bound up to an $H^{s-1}$
remainder contradicts the unit $H^s$ norms.
The Cauchy derivative varies with initial data
in operator norm uniformly on the fixed compact
time interval, so the nonlinear neighborhood
can also be chosen uniformly on a smaller
interval of delays.
\end{proof}

The orbit-only two-slice map still has the irrational
weak modes of Proposition~\ref{prop:both-ray-holonomy}
even on the circular backgrounds in this proposition.
For instance, the two projected directions
for a short circular loop are just the tangent
of the orbit curve and its rotated copy.
This makes the effect of the metric gradient visible
in the same evolution and at the same observation times.

\paragraph{Longer placements can make both full metrics weak.}
The shortness hypothesis has geometric content even
in the circular family. Use the smooth rotating
frame defined by parallel transport along the signed
radial geodesics, including at zero radius.
For any winding, choose the later time so that
\[
  \int_a^c k_0\cosh R(t)\,dt=\pi .
\]
For every small positive initial radius this defines
one delay near $\pi/k_0$, by strict monotonicity
in $c$. It is strictly shorter than $\pi/k_0$ and
depends analytically on the initial radius, so the
delays fill a nontrivial interval and are irrational
relative to $\pi$ outside a countable set of small
radii. To see that the endpoint loop can stay
nontrivial, take $a$ large with $k_0$ fixed.
The normalized linear radial solution with zero
initial derivative then approaches
$\cos(k_0(t-a))$ uniformly on this interval
as $a\to\infty$, by the scalar Duhamel formula
with damping $1/t$. Its terminal value is negative
for sufficiently large $a$ and remains so after
dividing the nearby nonlinear radius by its
initial amplitude.

At the chosen delay $W_d$ is the identity in
the rotating frame and the metric directions
at the two times are both angular, since the
single-ray rotation is a half turn. All their
translates are collinear. For the irrational
choices the two-full-metric criterion therefore
fails. This gives arbitrarily weak energy-normalized
velocity directions at longer exact placements,
even when both full induced metrics are supplied.
